% Canonical statement: sections/06_explicit.tex:11-22
\begin{lemma}[Small subgrid counting and decision cost]\label{lem:subgrid}
At grid scale $k\ge2$, put $K_k=2k^3$ and
\[
 I_k=2k^4-\left(\frac{k(k+1)}2\right)^2,
 \qquad D_k=\left\lfloor\frac{k}{24\log_2 k}\right\rfloor.
\]
Under independent fair flips,
\[
 \Pr[\sr(A)\le D_k]\le2^{-k^4/2}.
\]
Whether $\sr(A)\le D_k$ can be decided after reading at most $I_k\le2k^4$ oracle bits, in time $2^{O(k^3D_k\log k)}=2^{O(k^4)}$ when $D_k\ge1$. When $D_k=0$, the event is empty and is decided immediately.
\end{lemma}

% Proof binding: appendices/explicit_proofs.tex:4-20
\begin{proof}[Proof of \cref{lem:subgrid}]
When $D_k\ge1$, we have $k\ge32$ and $1\le D_k\le k\le K_k/2$, so the range condition in \eqref{eq:warren} holds. We use $k\ge32$ only to bound the logarithm, not as the first positive-floor threshold. Thus
\[
 \log_2(4eK_k/D_k)\le\log_2(8e k^3)\le4\log_2k.
\]
The strict-pattern bound \eqref{eq:warren} therefore gives at most
\[
 2^{2K_kD_k\cdot4\log_2k}\le2^{(2/3)k^4}
\]
low-rank tables. For $k\ge3$, $(k+1)^2/(4k^2)\le4/9<1/2$, so $I_k\ge(3/2)k^4$. Dividing the preceding count by $2^{I_k}$ gives at most $2^{-5k^4/6}\le2^{-k^4/2}$. If $D_k=0$, no nonempty strict sign matrix has that rank.

For the decision algorithm, query the incidence bits and fill the known off-incidence entries. The condition is equivalent to feasibility of
\[
 A_{ij}\sum_{s=1}^{D_k}U_{is}V_{sj}\ge1\quad(i,j\in[K_k]).
\]
Indeed a strict factorization can be scaled to have all signed scores at least one, since there are finitely many positive signed scores. The converse is immediate. There are $2K_kD_k$ real variables and $K_k^2$ degree-two polynomial inequalities with bounded integer coefficients. We use the singly exponential existential-real consistency bound in the formulation of \citet{Vorobjov21}, attributed there to \citet{Renegar92}: $(sd)^{O(v)}L^{O(1)}$ bit operations for $s$ polynomials of degree at most $d$, $v$ variables, and coefficient bit length $L$. Substitution gives $2^{O(k^3D_k\log k)}$. This is a use of the classical decision theorem, not a practical algorithm for the large full table.
\end{proof}
